\subsection{Higher-power Frobenius lines approaching both ends of the low-rate window}
\label{sec:higher-power-frobenius}

Full power-map fibers allow the message dimension to grow while the
source agreement approaches the first-order curve and the tested
agreement approaches Johnson. The alphabet below is quartic over its
prime field, and the evaluation domain is specified by multiplicative branches inside it.
The rate still vanishes. For a word $w$, write $A(w)$ for its maximum
agreement with the stated Reed--Solomon code, and write
$\mathcal L_T(w)$ for its list of degree-$<k$ polynomials with at least
$T$ matches. Ordinary common agreement is the largest number of
coordinates on which two words simultaneously agree with two such
polynomials.

\begin{lemma}[Noncanonical witnesses on odd-power blocks]
\label{lem:odd-power-exact-source}
Let $p\ge5$ be prime and let $h>1$ be an odd divisor of $M=p^2+1$
with $h<M$ and $h-1<M/h$. Put $B=\mathbb F_{p^2}\subset E=\mathbb F_{p^4}$,
choose primitive $\xi\in E^*$, and set
$D_0=\{x\in E^*:x^h\in B^*\}$ and $D_1=\xi D_0$.
On these blocks respectively put $(f,g)=(X^{hp},0)$ and
$(f,g)=(\xi^{h-hp}X^{hp},1)$. The following bound also holds after
arbitrary puncturing. If $\ell$ is the smallest prime divisor of $h$,
every polynomial $q$
of degree at most $h$ other than $aX^h+b$ satisfies
\[
 \operatorname{agr}(f+\lambda g,q)
 \le \frac h\ell\max\{p,h^2\}+(2h-h/\ell)h.
\]
In particular, if $p\ge h^2$, this bound is strictly less than $hp$.
\end{lemma}
\begin{proof}
We first bound a coefficientwise-$\mathbb F_{p^2}$ branch. Write
$B=\mathbb F_{p^2}$ and consider $z^{hp}=P(z)$, where $P\in B[Z]$
has degree at most $h$ and is not $aZ^h+b$.
Let $d$ be the greatest common divisor of $h$ and all root
multiplicities of $P$ over $\overline B$, and put $H=h/d$.
Then $P=R^d$ for a polynomial $R$ of degree at most $H$.
Since $\gcd(h,p^2-1)=1$, choose the unique $d$th root in $B$ of the
leading coefficient of $P$. With that leading coefficient, $R$ is fixed
by coefficientwise $p^2$-Frobenius, hence lies in $B[Z]$.
The original matching equation is equivalent on $B$ to $z^{Hp}=R(z)$.

The polynomial $U^H-R(Z)$ is irreducible over $\overline B(Z)$.
Indeed, $\overline B(Z)$ contains $\mu_H$; adjoining a root $u$ gives
a cyclic Galois extension, say of degree $e\mid H$.
If $e<H$, then $u^e$ belongs to the base field, so $R$ is a proper
power of order dividing $H$, contrary to the definition of $d$.
This also proves irreducibility in $\overline B[Z,U]$.
Writing $\sigma(c)=c^p$, all matches lie on
\[
 U^H=R(Z),\qquad Z^H=R^\sigma(U).
\]
If these degree-$H$ curves have no common component, B\'ezout gives
at most $H^2$ matches. Otherwise irreducibility and their equal degree
force their defining polynomials to be proportional, so
$R=aZ^H+b$. Since $H$th powering permutes $B$, the substitution
$y=z^H$ reduces the matching equation to $y^p=ay+b$, with at most
$p$ solutions. Thus each such branch contributes at most
$\max\{p,h^2\}$.

A noncanonical $q$ has a nonzero coefficient of some $X^j$ with
$0<j<h$. Within one block, this coefficient can be
coefficientwise in $B$ on at most $\gcd(j,h)\le h/\ell$ branches.
Write $\alpha_t=\xi^{(M/h)t}$. Good branches in both blocks would give
$q_j\alpha_t^j\in B$ and $q_j\xi^{j-h}\alpha_u^j\in B$, implying
$j-h+(M/h)j(u-t)\equiv0\pmod M$. This is impossible because
$0<|j-h|<M/h$.
All remaining branches contribute at most $h$ matches by projection
onto $E/B$. This proves the displayed bound. If $p\ge h^2$, dividing
it by $hp$ gives at most
$1/\ell+(2-1/\ell)/h<1$, since $h\ge\ell\ge3$.
\end{proof}


\begin{proposition}[Finite full-fiber criterion]
\label{prop:higher-power-frobenius-finite}
Let $p\ge5$ be prime, let $B=\mathbb F_{p^2}\subset E=\mathbb F_{p^4}$,
and let $h>1$ be an odd proper divisor of $M=p^2+1$ satisfying
$h-1<M/h$. Put $N=h(p^2-1)$. For an integer $0<m<N$, define
\[
 \begin{gathered}
 n=N+m,\qquad k=h+1,\qquad T=\lfloor\sqrt{hn-1}\rfloor,\\
 R=\left\lfloor\frac{m}{p^2-1}\right\rfloor,\qquad
 \tau=m-R(p^2-1).
 \end{gathered}
\]
Set $U=hp$ if $p\ge h^2$, and $U=hp+h^2$ otherwise. Suppose that
\begin{equation}\label{eq:higher-power-finite-guards}
 U<T\le(h+R)(p-1).
\end{equation}
There are an explicitly specified $n$-point domain $D\subset E$ and words $r_0,r_1\in E^D$
for which the dimension-$k$ Reed--Solomon code has
\[
 hp\le A(r_i)\le U\quad(i=0,1),\qquad
 A_{\mathrm{common}}(r_0,r_1)=hp.
\]
On the affine line $(1-z)r_0+zr_1$, exactly $(p+1)(p^2-1)$ parameters
$z\in E$ have singleton threshold-$T$ lists, one further parameter has
a list of size $p+1$, and all other parameters have empty lists.
Both endpoints have empty threshold-$T$ lists. If $p\ge h^2$, every
word with an empty threshold list has exact agreement $hp$. Moreover, $T$ is the
largest integer strictly below $\sqrt{(k-1)n}$. If $k/n$ is in the
low-rate branch of $a_1$ and
\begin{equation}\label{eq:higher-power-first-order-guard}
 \sqrt{kn/2}+(k^3n/72)^{1/4}<hp,
\end{equation}
then $n a_1(k/n)<hp$.
\end{proposition}

\begin{proof}
Choose a primitive $\xi\in E^*$, put $\eta=\xi^h$, and set
\[
 D_0=\{x\in E^*:x^h\in B^*\},\qquad D_1=\xi D_0.
\]
The quotient $E^*/B^*$ has order $M$. Since $h\mid M$, every
$y\in B^*$ has exactly $h$ preimages under $x\mapsto x^h$ in $E$.
Thus $|D_0|=|D_1|=N$. The two blocks are disjoint, and $\eta\notin B$,
because $h<M$. They decompose into the $B$-lines with representatives
$\alpha_t=\xi^{(M/h)t}$ and $\xi\alpha_t$, respectively, for
$0\le t<h$. Put $z_0=\xi^M$, a generator of $B^*$, and retain
\[
 D'_1=\bigcup_{0\le j<R}\xi\alpha_jB^*
 \;\cup\;\{\xi\alpha_Rz_0^v:0\le v<\tau\}.
\]
These disjoint sets have total size $m$. Thus the domain is determined
by $\xi$ and the given integers, without a probabilistic choice.
On $D=D_0\cup D'_1$ define
\[
 f(x)=
 \begin{cases}x^{hp},&x\in D_0,\\
 \eta^{1-p}x^{hp},&x\in D'_1,
 \end{cases}
 \qquad
 g(x)=\begin{cases}0,&x\in D_0,\\1,&x\in D'_1.\end{cases}
\]

For $a\in B$ with $a^{p+1}=1$, let
$I_a=\operatorname{Im}(y\mapsto y^p-ay)$ on $B$. This map has a
$p$-element kernel, and
\[
 I_a=\{v\in B:v^p=-a^pv\}.
\]
Consequently the $I_a$ are precisely the $p+1$ distinct
$\mathbb F_p$-lines in $B$. The polynomial $q_{a,b}=aX^h+b$ has
$hp$ matches with $f+\lambda g$ on $D_0$ when $b\in I_a\setminus\{0\}$,
and $h(p-1)$ when $b=0$. Writing $x^h=\eta y$ on $D_1$, its second
equation is $y^p-ay=(b-\lambda)/\eta$. Thus it has a full affine
fiber on each block exactly for
\begin{equation}\label{eq:higher-power-canonical-labels}
 a^{p+1}=1,\qquad b,v\in I_a,\qquad \lambda=b-\eta v.
\end{equation}
Since $h$ is odd and divides $p^2+1$, we have
$\gcd(h,p^2-1)=1$. On each complete branch $\xi\alpha_jB^*$,
the value $x^h/\eta$ therefore runs bijectively over $B^*$.
Each of the $R$ complete retained branches contributes $p$ matches
if $v\ne0$, and $p-1$ if $v=0$. The partial branch can only add matches.
Every polynomial in~\eqref{eq:higher-power-canonical-labels}
thus has at least $(h+R)(p-1)\ge T$ matches.

We next bound all other degree-$\le h$ polynomials. A polynomial
$aX^h+b$ not satisfying~\eqref{eq:higher-power-canonical-labels} has
at most $hp$ matches in total. Indeed, if $a\notin B$, two distinct
solutions of either equation in $y\in B$ would force $a\in B$;
if $a\in B$ and $a^{p+1}\ne1$, the map $y^p-ay$ is invertible.
These cases have at most $2h$ matches. In the remaining case each
block contributes one full affine fiber or no matches, so failure
of~\eqref{eq:higher-power-canonical-labels} leaves at most one fiber.

Now let $q(X)=\sum q_jX^j$ have some $q_j\ne0$ with $1\le j<h$.
On a core branch $\alpha_t B^*$, normalizing the matching equation
by $\alpha_t^{hp}\in B^*$ gives received monomial $z^{hp}$;
its coefficient of $z^j$ lies in $B$ only if $q_j\alpha_t^j\in B$.
On a fresh branch $\xi\alpha_u B^*$, normalize instead by
$\eta\alpha_u^{hp}$; the corresponding condition is
$q_j\xi^{j-h}\alpha_u^j\in B$. If both conditions held, then
\[
 j-h+(M/h)j(u-t)\equiv0\pmod M.
\]
Solvability requires $(M/h)\gcd(j,h)\mid j-h$, impossible because
$0<h-j\le h-1<M/h$. Hence one whole block has no branch whose
normalized matching polynomial has all coefficients in $B$.
On each branch of that block, a nonzero $B$-linear map $E\to B$
with kernel $B$ kills the received monomial and leaves a nonzero
polynomial of degree at most $h$. This gives at most $h^2$ matches
on that block. The other block contributes at most $hp$ by the
ordinary degree bound. Puncturing preserves these bounds, proving
that every such $q$ has at most $hp+h^2$ matches. When $p\ge h^2$,
Lemma~\ref{lem:odd-power-exact-source} sharpens this to fewer than $hp$
matches. In either case the bound is at most $U<T$.

Since $E=B\oplus\eta B$, the planes $I_a+\eta I_a$ intersect pairwise
only at zero. Each nonzero label in their union therefore determines
exactly one pair $(a,b)$ in~\eqref{eq:higher-power-canonical-labels}.
There are $(p+1)(p^2-1)$ such labels. At zero the witnesses are the
$p+1$ polynomials $aX^h$, and outside this union the list is empty.
For explicit endpoints, take $\lambda_0=1+\eta z_0$ and
$\lambda_1=2+\eta z_0$, and put $r_i=f+\lambda_i g$.
Since $z_0\notin\mathbb F_p$, neither label lies in a plane
$I_a+\eta I_a$: its two $B$-components are linearly independent
over $\mathbb F_p$. The preceding bounds give $A(r_i)\le U$,
and a core fiber with $b\ne0$ gives $A(r_i)\ge hp$.
The same argument gives exact agreement $hp$ at every nonexceptional
parameter whenever $p\ge h^2$.

To bound ordinary common agreement for $(f,g)$, a nonconstant
degree-$\le h$ polynomial explaining $g$ has at most $h$ zero-values
on $D_0$ and $h$ one-values on $D'_1$. A constant explanation confines
the common support to one block, where explaining $f$ has at most
$hp$ solutions. The core fiber with $b\ne0$, together with the zero
explanation of $g$, attains $hp$. Invertible linear changes of the
two received words preserve ordinary common agreement, while
$\lambda=\lambda_0+z(\lambda_1-\lambda_0)$ is a bijection on $E$.
This proves the asserted profile and common agreement for $r_0,r_1$.
Finally, the low-rate formula $a_1(\rho)=t(1+u)$, where
$t=\sqrt{\rho/2}$ and $u^2(u+3)=t$, gives $u\le\sqrt{t/3}$ and hence
\[
 n a_1(k/n)\le\sqrt{kn/2}+(k^3n/72)^{1/4}.
\]
\end{proof}

\begin{theorem}\label{thm:higher-power-frobenius-window}
There is a sequence of Reed--Solomon codes over $E=\mathbb F_{p^4}$,
on explicitly specified domains of size $n$, with growing dimensions $k=n^{o(1)}$
and rate $k/n=n^{-1+o(1)}$, together with received words $r_0,r_1$ and
integer thresholds $T$, such that
\[
 n a_1(k/n)<A_{\mathrm{common}}(r_0,r_1)
 =A(r_0)=A(r_1)<T<\sqrt{(k-1)n}.
\]
The affine line between these words has $n^{3/2-o(1)}$ singleton
threshold-$T$ lists, one further list of size $p+1$, and empty lists
elsewhere; every word in the latter class has exact agreement $hp$.
Here $h=k-1$. Moreover, for $i=0,1$,
\[
 \frac{A(r_i)}{n a_1(k/n)}\longrightarrow1,\qquad
 \frac{A_{\mathrm{common}}(r_0,r_1)}{n a_1(k/n)}\longrightarrow1,
 \qquad \frac{T}{\sqrt{(k-1)n}}\longrightarrow1,
\]
\[
 \frac{T-A(r_i)}{T-k}\longrightarrow1-\frac1{\sqrt2},\qquad
 \frac{T-A_{\mathrm{common}}(r_0,r_1)}{T-k}
 \longrightarrow1-\frac1{\sqrt2}.
\]
The alphabet size is $|E|=n^{2-o(1)}$. These are limits for the
locations of the source and tested agreements, not an assertion of
optimality for a uniform exceptional-parameter count.
\end{theorem}

\begin{proof}
Take $h=5^a$ with $a\to\infty$. A square root of $-1$ modulo $5$
lifts to one modulo $h$. Dirichlet's theorem allows an increasing
sequence of primes in this residue class with $p>h^{a+4}$. Thus
$h\mid p^2+1$, $h=p^{o(1)}$, and $p/h^4\to\infty$.
Set $c=1-2/\sqrt h$ and $m=\lfloor cN\rfloor$ in
Proposition~\ref{prop:higher-power-frobenius-finite}. Its parameters obey
\[
 n=(2-2/\sqrt h+o(1))hp^2,\qquad
 T=(\sqrt2+o(1))hp,\qquad R/h\longrightarrow1.
\]
Both inequalities in~\eqref{eq:higher-power-finite-guards} hold eventually:
$p\ge h^2$, so $U=hp$, and $(h+R)(p-1)\sim2hp$. The branch condition also holds.

For first-order placement, $n\le(2-2/\sqrt h)hp^2$, and for $h\ge25$
\[
 \sqrt{kn/2}
 \le hp\sqrt{(1+1/h)(1-1/\sqrt h)}
 \le hp-\frac25p\sqrt h.
\]
The correction $(k^3n/72)^{1/4}=O(h\sqrt p)$ is smaller than this
deficit. Thus~\eqref{eq:higher-power-first-order-guard} holds.
The lower bound $a_1(\rho)>\sqrt{\rho/2}$ and the same upper bound
show $n a_1(k/n)\sim hp$. The finite proposition gives
$A_{\mathrm{common}}=A(r_i)=hp$, proving the claimed
agreement ratios. Also $k=o(hp)$, so both loss ratios tend to
$1-1/\sqrt2$. Finally $n=p^{2+o(1)}$ and the singleton count is
$(p+1)(p^2-1)$, yielding all stated exponents.
\end{proof}

\paragraph{Exact finite parameter certificates.}
A smaller instance has $p=97$, $h=5$, $m=23520$, $n=70560$,
$k=6$, and $T=593$. Its common and both individual endpoint agreements are exactly $485$,
and there are $921984$ singleton
parameters. The integer first-order upper bound is $483$, while the
minimum canonical agreement is $672$. Thus the finite inequalities
already hold at these parameters.


For the BabyBear prime $p=2013265921$, take $h=12241$ and
$m=\lfloor54h(p^2-1)/55\rfloor$. Exact integer checks of the finite
criterion give
\[
 \begin{aligned}
 n&=98329309808423281932846,& k&=12242,\\
 A_{\mathrm{common}}&=24644388138961,&
 A(r_i)&=24644388138961,\\
 T&=34693646123820,&
 \#\{\text{singleton labels}\}&=8160249298611705595853537280.
 \end{aligned}
\]
The first-order upper bound is $24533338196068<A_{\mathrm{common}}$,
and the deterministic retention rule has $R=12018$ complete branches,
$\tau=1768686400869808686$ further points, and minimum canonical
agreement $48839817953280>T$. The guaranteed endpoint loss divided
by the capacity margin is
\[
 \frac{T-hp}{T-k}
 =\frac{10049257984859}{34693646111578}>0.289657.
\]
The branch-prefix formula specifies the retained domain directly; this
parameter certificate does not enumerate its coordinates. Its rate is
approximately $1.25\cdot10^{-19}$. More generally, the asymptotic
construction has capacity margin $(T-k)/n=\Theta(p^{-1})$ and
fractional proximity loss $\Theta(p^{-1})$, despite a limiting relative
loss of $1-1/\sqrt2\approx29.3\%$. It does not give a constant-rate,
prime-alphabet, or prescribed short-domain construction.
